\section{第三章：线性方程}
\basis{第3.1--3.5节，教材第60--88页。}
\problem{3.1.1}{64--65}{证明可微矩阵函数 $A(t)=(a_{ij}(t))$ 的行列式导数等于依次将一行换成该行导数所得的 $n$ 个行列式之和。}
\begin{proof}
由 Leibniz 展开式
\[
 \det A(t)=\sum_{\sigma\in S_n}\operatorname{sgn}(\sigma)
 \prod_{j=1}^na_{j,\sigma(j)}(t).
\]
对每个有限乘积求导：
\[
 (\det A)'=\sum_{k=1}^n\sum_{\sigma\in S_n}\operatorname{sgn}(\sigma)
 a'_{k,\sigma(k)}\prod_{j\ne k}a_{j,\sigma(j)}.
\]
固定 $k$ 后内层和正是第 $k$ 行被其导数替换后的行列式，逐项相加即为所求公式。该式不要求 $A$ 可逆，也不等于 $\det A'$。
这里所有求和都是有限和，因此可以先对每个排列项使用乘积求导，再交换“选取求导行 $k$”与“遍历排列 $\sigma$”的顺序。固定 $k$ 后，第 $k$ 行的元素全部换为相应导数，其余行保持原样，排列符号仍为 $\operatorname{sgn}(\sigma)$，所以内层和完整地对应一个行列式，而没有遗漏交叉项。
\end{proof}
\problem{3.1.2}{65}{设 $A=\begin{pmatrix}2t^3+1&t^2\\t&4t^2\end{pmatrix}$，计算并比较 $\det A'$ 与 $(\det A)'$。}
先对矩阵的每个元素求导：
\[
A'=\begin{pmatrix}6t^2&2t\\1&8t\end{pmatrix}.
\]
“先求导再取行列式”给出
\[
\det A'=(6t^2)(8t)-(2t)\cdot1=48t^3-2t.
\]
“先取行列式再求导”则分两步：
\[
\begin{aligned}
\det A&=(2t^3+1)(4t^2)-t^2\cdot t
=8t^5+4t^2-t^3,\\
(\det A)'&=40t^4+8t-3t^2.
\end{aligned}
\]
两种运算产生的各项次数不同，不能将它们交换。
\ans{\det A'=48t^3-2t,\qquad (\det A)'=40t^4-3t^2+8t.}
例如 $t=1$ 时两者分别为46与45，通常不相等。
\problem{3.1.3}{65}{设 $x$ 在 $[\alpha,\beta]$ 连续，且 $|x(t)|\le L+M\displaystyle\int_\alpha^t|x(s)|\dd s$，$L,M\ge0$。用逐次逼近证明 $|x(t)|\le L\e^{M(t-\alpha)}$。}
\begin{proof}
记 $K=\max|x(t)|$。反复将原不等式代入右端积分，归纳得到
\[
 |x(t)|\le L\sum_{j=0}^{n-1}\dfrac{[M(t-\alpha)]^j}{j!}
 +K\dfrac{[M(t-\alpha)]^n}{n!}.
\]
归纳一步使用 $\displaystyle\int_\alpha^t(s-\alpha)^j\dd s=(t-\alpha)^{j+1}/(j+1)$。由于有限区间上阶乘增长快于固定底数幂，最后一项一致趋零，前一项趋于指数级数，故得结论。
具体前两次代入为
\[
|x(t)|\le L+KM(t-\alpha),
\]
\[
|x(t)|\le L+LM(t-\alpha)+\dfrac{KM^2(t-\alpha)^2}{2}.
\]
再代入一次就增加 $L[M(t-\alpha)]^2/2!$，并把余项次数提高到3；这解释了归纳式的每一个系数从何而来。
\end{proof}
\problem{3.2.1}{70--71}{验证 $X(t)=\begin{pmatrix}\e^{2t}&t\e^{2t}\\0&\e^{2t}\end{pmatrix}$ 是 $x'=\begin{pmatrix}2&1\\0&2\end{pmatrix}x$ 的基本解矩阵。}
基本解矩阵必须同时满足 $X'=AX$ 与 $\det X\ne0$，故分开检查。第一步逐项求导，尤其对 $t\e^{2t}$ 使用乘积法则：
\[
X'=\begin{pmatrix}2\e^{2t}&\e^{2t}+2t\e^{2t}\\0&2\e^{2t}\end{pmatrix}.
\]
第二步直接做矩阵乘法：
\[
AX=\begin{pmatrix}2&1\\0&2\end{pmatrix}
\begin{pmatrix}\e^{2t}&t\e^{2t}\\0&\e^{2t}\end{pmatrix}
=\begin{pmatrix}2\e^{2t}&(2t+1)\e^{2t}\\0&2\e^{2t}\end{pmatrix}=X'.
\]
最后算 $\det X=\e^{2t}\e^{2t}-0=\e^{4t}>0$。因此两列是线性无关的解；代入 $t=0$ 可见 $X(0)=I$。
\ans{x=X(t)c,\qquad X(0)=I.}
两列既为解又线性无关，所以是基本解矩阵，且在0处标准化。
\problem{3.2.2}{71}{两个齐次方程组 $x'=A_1(t)x$、$x'=A_2(t)x$ 有同一个基本解矩阵，证明 $A_1=A_2$。}
\begin{proof}
设共同基本解矩阵为 $X$，则 $A_1X=X'=A_2X$。每个 $t$ 上 $X$ 可逆，右乘 $X^{-1}$ 即得 $A_1=A_2$。
把最后一步完整写为 $(A_1-A_2)X=0$，右乘 $X^{-1}$ 得 $(A_1-A_2)(XX^{-1})=0X^{-1}$，即 $A_1-A_2=0$。这里基本解矩阵的可逆性不可少；仅有一个共同的向量解不足以消去 $X$。
\end{proof}
\problem{3.2.3}{71}{若 $(\sin t,\cos t)^T,(\cos t,-\sin t)^T$ 是 $x'=(a_{ij}(t))x$ 的基本解组，求各 $a_{ij}$。}
把两个已知列解组成 $X$，由 $X'=AX$ 可反求 $A=X'X^{-1}$。先计算
\[
\det X=-\sin^2t-\cos^2t=-1,
\qquad X^{-1}=\begin{pmatrix}\sin t&\cos t\\\cos t&-\sin t\end{pmatrix}=X.
\]
再算 $X'=\begin{pmatrix}\cos t&-\sin t\\-\sin t&-\cos t\end{pmatrix}$。逐个计算乘积的四个元素：
\[
\begin{aligned}
a_{11}&=\cos t\sin t-\sin t\cos t=0,\\
a_{12}&=\cos^2t+\sin^2t=1,\\
a_{21}&=-\sin^2t-\cos^2t=-1,\\
a_{22}&=-\sin t\cos t+\cos t\sin t=0.
\end{aligned}
\]
\ans{A=\begin{pmatrix}0&1\\-1&0\end{pmatrix}.}
\checknote{两列分别满足 $x_1'=x_2,x_2'=-x_1$；$\det X=-1$，反演合法。}
\problem{3.2.4}{71}{用存在唯一性证明：$x'=A(t)x$ 的 $n$ 个解线性无关，当且仅当其 Wronski 行列式在某一点不为零。}
\begin{proof}
若 $\det X(t_0)\ne0$，而 $\sum c_kx_k(t)=0$，代入 $t_0$ 得 $X(t_0)c=0$，所以 $c=0$，解组线性无关。

反之，若 $\det X(t_0)=0$，存在 $c\ne0$ 使 $X(t_0)c=0$。线性组合 $z=Xc$ 满足 $z'=Az,z(t_0)=0$，由唯一性 $z\equiv0$，故解组线性相关。于是线性无关时行列式不但在某一点，而且在每一点都非零。
在反向论证中，$c$ 必须是由单一时刻选出的常向量。正因 $c$ 不随 $t$ 变化，$z'=\sum c_kx_k'=A\sum c_kx_k=Az$ 才成立；不能让不同时间使用不同的线性相关系数来代替这一步。
\end{proof}
\problem{3.2.6}{71}{$A(t)$ 是 $\omega$ 周期连续矩阵，$X(t)$ 为 $x'=A(t)x$ 的基本解矩阵。证明 $X(t+\omega)$ 也是基本解矩阵，且存在可逆常矩阵 $C$ 使 $X(t+\omega)=X(t)C$。}
\begin{proof}
周期性给出
\[
 \dfrac\dd{\dd t}X(t+\omega)=A(t+\omega)X(t+\omega)=A(t)X(t+\omega).
\]
其行列式非零，所以是基本解矩阵。令 $C(t)=X^{-1}(t)X(t+\omega)$，使用 $(X^{-1})'=-X^{-1}A$，得到 $C'=0$，因此 $C$ 为常矩阵。两因子均可逆，故 $C$ 可逆，且
\[
 C=X^{-1}(t_0)X(t_0+\omega).
\]
所用的乘积求导可逐项写成
\[
\begin{aligned}
C'&=(X^{-1})'X(t+\omega)+X^{-1}X'(t+\omega)\\
&=-X^{-1}AX(t+\omega)+X^{-1}AX(t+\omega)=0.
\end{aligned}
\]
因此固定任意 $t_0$ 后，所有时刻都使用同一 $C=X^{-1}(t_0)X(t_0+\omega)$。
\end{proof}
\problem{3.3.1}{73--74}{设 $x'=A(t)x+f(t)$，其中
\[
 A=\begin{pmatrix}\cos^2t&\dfrac12\sin2t-1\\\dfrac12\sin2t+1&\sin^2t\end{pmatrix},\quad
 f=\binom{\cos t}{\sin t}.
\]
(1) 验证 $\Phi=\begin{pmatrix}\e^t\cos t&-\sin t\\\e^t\sin t&\cos t\end{pmatrix}$ 为齐次方程组的基本解矩阵；(2) 求 $x(0)=(-1,2)^T$ 的解。}
\begin{enumerate}
\item 令 $r=(\cos t,\sin t)^T,s=(-\sin t,\cos t)^T$，于是 $r'=s,s'=-r$。利用 $\dfrac12\sin2t=\sin t\cos t$，逐行相乘：
\[
\begin{aligned}
(Ar)_1&=\cos^3t+\sin^2t\cos t-\sin t=\cos t-\sin t,\\
(Ar)_2&=\sin t\cos^2t+\cos t+\sin^3t=\sin t+\cos t,\\
(As)_1&=-\cos^2t\sin t+\sin t\cos^2t-\cos t=-\cos t,\\
(As)_2&=-\sin^2t\cos t-\sin t+\sin^2t\cos t=-\sin t.
\end{aligned}
\]
所以 $Ar=r+s,As=-r$。由此 $(\e^tr)'=\e^t(r+s)=A(\e^tr)$，且 $s'=As$，两列均是解。再算 $\det\Phi=\e^t(\cos^2t+\sin^2t)=\e^t\ne0$，故是基本解矩阵。
\item 令 $x=\Phi c(t)$，代入非齐次方程并约去 $\Phi'c=A\Phi c$，得 $\Phi c'=f$。本题
\[
\Phi^{-1}=\begin{pmatrix}\e^{-t}\cos t&\e^{-t}\sin t\\-\sin t&\cos t\end{pmatrix},
\quad c'=\Phi^{-1}\binom{\cos t}{\sin t}=\binom{\e^{-t}}0.
\]
由于 $\Phi(0)=I$，有 $c(0)=(-1,2)^T$。采用定积分计算：
\[
c_1(t)=-1+\int_0^t\e^{-s}\dd s=-\e^{-t},\qquad c_2(t)=2.
\]
最后乘回 $\Phi$，第一分量为 $\e^t\cos t(-\e^{-t})-2\sin t$，第二分量为 $\e^t\sin t(-\e^{-t})+2\cos t$，化简如下。
\ans{x(t)=\binom{-\cos t-2\sin t}{-\sin t+2\cos t}.}
\end{enumerate}
\checknote{$t=0$ 得 $(-1,2)^T$；用 $Ar=r+s,As=-r$ 可核验 $(-r+2s)'=A(-r+2s)+r$。}
\problem{3.3.2}{74}{$A(t),f(t)$ 在 $[\alpha,\beta]$ 连续，且 $f$ 不恒为零。证明非齐次组 $x'=Ax+f$ 存在且至多存在 $n+1$ 个线性无关的解。}
\begin{proof}
取一个特解 $x_p$ 和齐次基本解组 $u_1,\ldots,u_n$。每个非齐次解均为 $x_p+\sum c_iu_i$，所以所有非齐次解都在至多 $n+1$ 维空间 $\operatorname{span}\{x_p,u_1,\ldots,u_n\}$ 中。

考虑 $x_p,x_p+u_1,\ldots,x_p+u_n$。若它们有线性关系，则
\[
 \left(a_0+\sum_{i=1}^na_i\right)x_p+\sum_{i=1}^na_iu_i=0.
\]
作用算子 $Lx=x'-Ax$，得到 $(a_0+\sum a_i)f=0$。由于 $f$ 不恒为零，括号为零，于是 $\sum a_iu_i=0$，从而所有 $a_i=0$，再得 $a_0=0$。故恰可构造 $n+1$ 个线性无关的解。
“$f$ 不恒为零”表示存在一个时刻 $t_*$，使 $f(t_*)\ne0$。将算子得到的向量恒等式在 $t_*$ 处取值，立刻给出常标量 $a_0+\sum a_i=0$；无须假设 $f(t)$ 在每一点均非零。随后用的是齐次基本解组的函数线性无关性。
\end{proof}
\problem{3.3.3}{74}{证明 $x'=A(t)x+f(t,x),x(t_0)=x_0$ 等价于
\[
 x(t)=X(t)X^{-1}(t_0)x_0+\int_{t_0}^tX(t)X^{-1}(s)f(s,x(s))\dd s,
\]
其中 $X$ 是对应齐次方程的基本解矩阵。}
\begin{proof}
若 $x$ 满足初值问题，令 $c=X^{-1}x$，则
\[
 c'=(X^{-1})'x+X^{-1}x'=X^{-1}f(t,x(t)).
\]
从 $t_0$ 到 $t$ 积分再左乘 $X(t)$，即得所列积分式。

反之，若连续函数 $x$ 满足该积分式，由被积函数连续，右端可微。求导时 $X'=AX$，上限求导项为 $f(t,x(t))$，故 $x'=Ax+f(t,x)$。令 $t=t_0$ 又得 $x(t_0)=x_0$，所以两者等价。
反向求导时，积分项的导数具体为
\[
A(t)\int_{t_0}^tX(t)X^{-1}(s)f(s,x(s))\dd s+f(t,x(t)).
\]
第一项是积分式中相应部分左乘 $A(t)$，上限项中 $X(t)X^{-1}(t)=I$。再与初值项的导数合并，恰为 $A(t)x(t)+f(t,x(t))$。
\end{proof}
\problem{3.4.1}{79}{若 $x_s$ 分别满足 $Lx_s=f_s$，$s=1,\ldots,m$，其中 $L=D^n+a_1(t)D^{n-1}+\cdots+a_n(t)$，证明 $\sum C_sx_s$ 满足 $Lx=\sum C_sf_s$。}
\begin{proof}
$C_s$ 为常数，有限次求导可以逐项进行。由每一项乘法与加法的分配律，
\[
 L\left(\sum_s C_sx_s\right)=\sum_s C_sLx_s=\sum_s C_sf_s,
\]
即得结论。这里若所有 $f_s$ 相同，则右端为 $(\sum C_s)f$，不能直接写成 $f$。
若把最高阶项也写出，这个等式来自
\[
\left(\sum_sC_sx_s\right)^{(k)}=\sum_sC_sx_s^{(k)},\qquad k=0,\ldots,n.
\]
将它逐项代入 $L$，先对有限个 $s$ 交换求和，再将每个括号识别为 $Lx_s=f_s$，没有对变系数 $a_j(t)$ 进行额外求导。
\end{proof}
\problem{3.4.2}{79}{若 $x$ 为 $x''+a_1(t)x'+a_2(t)x=0$ 的非零解，证明 $x(t_0)=0$ 时 $x'(t_0)\ne0$。}
\begin{proof}
若同时 $x'(t_0)=0$，把方程化为 $(x,x')$ 的一阶线性方程组。零函数满足相同的两个初值，由唯一性，$x\equiv0$，与“非零解”矛盾。故 $x'(t_0)\ne0$。
使用唯一性所对应的系统是
\[
\binom{x}{x'}'=\begin{pmatrix}0&1\\-a_2(t)&-a_1(t)\end{pmatrix}\binom{x}{x'}.
\]
若两个初值都为0，该系统的初始向量就是零向量；零解满足相同系统与初值，因此在整个连通存在区间内与原解相同。
\end{proof}
\problem{3.4.3}{79}{验证 $x=\sin t/t$ 是 $x''+2x'/t+x=0$ 的解，并求通解。}
由于已知解含 $1/t$，在 $t\ne0$ 处令 $x=u/t$，逐次求导：
\[
\begin{aligned}
x'&=\dfrac{u'}t-\dfrac u{t^2},\\
x''&=\dfrac{u''}t-\dfrac{2u'}{t^2}+\dfrac{2u}{t^3}.
\end{aligned}
\]
把各项完整代入并抵消：
\[
\begin{aligned}
x''+\dfrac2t x'+x
&=\dfrac{u''}t-\dfrac{2u'}{t^2}+\dfrac{2u}{t^3}
+\dfrac{2u'}{t^2}-\dfrac{2u}{t^3}+\dfrac ut\\
&=\dfrac{u''+u}{t}.
\end{aligned}
\]
$u=\sin t$ 满足 $u''+u=0$，验证了给定解。该简谐方程的全部解为 $u=C_1\cos t+C_2\sin t$；乘回 $1/t$ 即得到全部 $x$。
\ans{x=\dfrac{C_1\cos t+C_2\sin t}{t},\quad t\ne0.}
两基解 Wronski 行列式为 $1/t^2$，在所取区间非零。
\problem{3.4.4}{79}{不使用 Liouville 公式，直接令 $x=x_1(t)y$，对二阶齐次方程降阶并求通解，其中 $x_1$ 为已知非零解。}
在 $x_1\ne0$ 的区间令 $x=x_1y$，先把两个导数完全展开：
\[
x'=x_1'y+x_1y',\qquad x''=x_1''y+2x_1'y'+x_1y''.
\]
代入后按照 $y,y',y''$ 分组：
\[
(x_1''+a_1x_1'+a_2x_1)y+(2x_1'+a_1x_1)y'+x_1y''=0.
\]
第一组为零，因为 $x_1$ 是已知解。令 $v=y'$，将剩下的式子除以 $x_1$，得到
\[
v'+\left(2\dfrac{x_1'}{x_1}+a_1\right)v=0.
\]
于是非零分支有
\[
\ln|v|=-2\ln|x_1|-\int a_1\dd t+C,
\qquad v=C_2\dfrac{\e^{-\int a_1\dd t}}{x_1^2}.
\]
再积分 $y'=v$ 时必须加入另一个独立常数 $C_1$，最后乘回 $x_1$；$C_2=0$ 包括比例于 $x_1$ 的原解。
\ans{x=x_1(t)\left[C_1+C_2\int\dfrac{\e^{-\int a_1(t)\dd t}}{x_1(t)^2}\dd t\right].}
第二基解与 $x_1$ 的 Wronski 行列式是非零常数倍的 $\displaystyle \e^{-\int a_1\dd t}$，故独立；$x_1$ 的孤立零点处用原方程延拓。
\problem{3.4.5}{79}{设 $x_1,x_2$ 是 $x''+a_1x'+a_2x=f$ 对应齐次方程的基本解组，证明常数变易给出的通解表达式。}
\begin{proof}
设 $W=x_1x_2'-x_1'x_2\ne0$。令 $x=c_1(t)x_1+c_2(t)x_2$，附加约束 $c_1'x_1+c_2'x_2=0$，使一阶导数简化为 $x'=c_1x_1'+c_2x_2'$。再求导并代入方程，利用两个齐次解满足的方程，剩下
\[
\begin{cases}
x_1c_1'+x_2c_2'=0,\\
x_1'c_1'+x_2'c_2'=f.
\end{cases}
\]
用二阶行列式解出两个未知导数：
\[
c_1'=\dfrac{\begin{vmatrix}0&x_2\\f&x_2'\end{vmatrix}}W=-\dfrac{x_2f}W,
\quad
c_2'=\dfrac{\begin{vmatrix}x_1&0\\x_1'&f\end{vmatrix}}W=\dfrac{x_1f}W.
\]
分别从 $t_0$ 到 $t$ 积分，得到 $\displaystyle c_1=C_1-\int_{t_0}^t x_2(s)f(s)/W(s)\dd s$、$\displaystyle c_2=C_2+\int_{t_0}^t x_1(s)f(s)/W(s)\dd s$。代回 $x$，将两个定积分合并为同一被积函数，得到下式。
\ans{x(t)=C_1x_1(t)+C_2x_2(t)+
\int_{t_0}^t\dfrac{x_1(s)x_2(t)-x_1(t)x_2(s)}{W(s)}f(s)\dd s.}
积分项及其一阶导数在 $t_0$ 均为零，而其二阶求导的上限项为 $f(t)$，故它为特解。任意另一解与此特解之差为齐次解，所以表达式涵盖全部解。
\end{proof}
\problem{3.4.6}{79--80}{已知基本解组 $\cos2t,\sin2t$，求 $x''+4x=t\sin2t$ 的通解。}
先算 Wronski 行列式：
\[
W=\begin{vmatrix}\cos2t&\sin2t\\-2\sin2t&2\cos2t\end{vmatrix}
=2\cos^22t+2\sin^22t=2.
\]
常数变易给 $c_1'=-t\sin^22t/2$、$c_2'=t\sin2t\cos2t/2$。用倍角公式后，逐个分部积分：
\[
\begin{aligned}
c_1&=-\dfrac14\int t(1-\cos4t)\dd t
=-\dfrac{t^2}{8}+\dfrac{t\sin4t}{16}+\dfrac{\cos4t}{64},\\
c_2&=\dfrac14\int t\sin4t\dd t
=-\dfrac{t\cos4t}{16}+\dfrac{\sin4t}{64}.
\end{aligned}
\]
两个积分常数暂取零，最后收入齐次通解。由 $x_p=c_1\cos2t+c_2\sin2t$，其含 $t$ 的混合项为
\[
\dfrac t{16}(\sin4t\cos2t-\cos4t\sin2t)=\dfrac t{16}\sin2t,
\]
其不含 $t$ 的混合项为 $\cos2t/64$，属于齐次解，可收入 $C_1\cos2t$。剩余特解为 $-t^2\cos2t/8+t\sin2t/16$。
\ans{x=C_1\cos2t+C_2\sin2t-\dfrac{t^2}{8}\cos2t+\dfrac t{16}\sin2t.}
\checknote{$(D^2+4)(t^2\cos2t)=2\cos2t-8t\sin2t$，$(D^2+4)(t\sin2t)=4\cos2t$，代入即为右端。}
\problem{3.5.1}{87}{求 Airy 方程 $x''-tx=0$ 在 $t=0$ 展开的幂级数解，$t\in\R$。}
设 $x=\sum_{n=0}^\infty a_nt^n$，逐项求导后把幂次统一成 $t^n$：
\[
x''=\sum_{n=0}^\infty(n+2)(n+1)a_{n+2}t^n,
\qquad tx=\sum_{n=1}^\infty a_{n-1}t^n.
\]
因此常数项为 $2a_2=0$；$n\ge1$ 的系数满足 $(n+2)(n+1)a_{n+2}=a_{n-1}$。令指标平移一次，写为
\[
a_{n+3}=\dfrac{a_n}{(n+3)(n+2)},\qquad n\ge0.
\]
分别沿三个余数类递推：
\[
\begin{aligned}
a_3&=a_0/6,&a_6&=a_3/30=a_0/180,&a_9&=a_6/72=a_0/12960,\\
a_4&=a_1/12,&a_7&=a_4/42=a_1/504,&a_{10}&=a_7/90=a_1/45360,\\
a_2&=0,&a_5&=0,&a_8&=0.
\end{aligned}
\]
所以 $a_0,a_1$ 独立决定两组级数，第三组全零；把相同自由系数的项分别收集即可。
\ans{x=a_0\left(1+\dfrac{t^3}{6}+\dfrac{t^6}{180}+\dfrac{t^9}{12960}+\cdots\right)
+a_1\left(t+\dfrac{t^4}{12}+\dfrac{t^7}{504}+\dfrac{t^{10}}{45360}+\cdots\right),\\
a_{3k}=\dfrac{a_0}{\prod_{j=1}^k(3j)(3j-1)},\qquad
 a_{3k+1}=\dfrac{a_1}{\prod_{j=1}^k(3j+1)(3j)}.}
空积取1。每个子级数相邻项绝对值之比对固定 $t$ 趋于零，故收敛半径为无穷；两解在0处初值为 $(1,0)$ 与 $(0,1)$，故独立。
\problem{3.5.2}{87}{用级数法求 $(t^2-2t)x''+5(t-1)x'+3x=0$，$x(1)=10,x'(1)=3$。}
展开中心应取初值点 $t=1$。令 $s=t-1$，有 $t^2-2t=s^2-1$，且对 $s$、$t$ 求导相同。设 $x=\sum_{n\ge0}a_ns^n$，则
\[
\begin{aligned}
x_{ss}&=\sum_{n\ge0}(n+2)(n+1)a_{n+2}s^n,\\
s^2x_{ss}&=\sum_{n\ge0}n(n-1)a_ns^n,\\
5sx_s&=\sum_{n\ge0}5na_ns^n.
\end{aligned}
\]
代入并比较同次幂，得到
\[
-(n+2)(n+1)a_{n+2}+[n(n-1)+5n+3]a_n=0.
\]
括号为 $(n+1)(n+3)$，约去非零的 $n+1$ 得 $a_{n+2}=(n+3)a_n/(n+2)$。初值给 $a_0=10,a_1=3$，因此
\[
a_2=\dfrac32\cdot10=15,\quad a_3=\dfrac43\cdot3=4,
\quad a_4=\dfrac54\cdot15=\dfrac{75}{4},\quad a_5=\dfrac65\cdot4=\dfrac{24}{5}.
\]
按奇偶两类一直递推，得到下列级数与乘积公式。
\ans{x=10+3s+15s^2+4s^3+\dfrac{75}{4}s^4+\dfrac{24}{5}s^5+\cdots,\quad s=t-1,\ |s|<1,\\
a_{2k}=10\prod_{j=1}^k\dfrac{2j+1}{2j},\qquad
 a_{2k+1}=3\prod_{j=1}^k\dfrac{2j+2}{2j+1}.}
相邻同奇偶项的系数比趋于1，且两类非零，故收敛半径恰为1；区间 $0<t<2$ 包含初值点，并避开首项系数零点。
\problem{3.5.3}{87}{解 Hermite 方程 $x''-2tx'+\lambda x=0$，并证明 $\lambda=2n$ 时存在次数为 $n$ 的多项式解。}
设 $x=\sum_{k\ge0}a_kt^k$，其中 $a_0=x(0),a_1=x'(0)$。三项在同次幂下分别贡献
\[
x'':(k+2)(k+1)a_{k+2},\qquad -2tx':-2ka_k,
\qquad \lambda x:\lambda a_k.
\]
故 $(k+2)(k+1)a_{k+2}+(\lambda-2k)a_k=0$，即
\[
a_{k+2}=\dfrac{2k-\lambda}{(k+2)(k+1)}a_k.
\]
前几项为 $a_2=-\lambda a_0/2$、$a_3=(2-\lambda)a_1/6$、$a_4=(4-\lambda)a_2/12$。偶数项全由 $a_0$ 产生，奇数项全由 $a_1$ 产生。固定 $t$ 时每类相邻项之比为 $O(|t|^2/k)$，故两级数处处收敛。

当 $\lambda=2n$，若 $n$ 偶则取 $a_0\ne0,a_1=0$；若 $n$ 奇则取 $a_0=0,a_1\ne0$。同奇偶的递推在 $k<n$ 时分子不为零，所以能达到 $a_n\ne0$；在 $k=n$ 时分子变零，所以 $a_{n+2}$ 及此后的同奇偶项全部为零。另一个奇偶类从初值起就全零，故所得解确为次数 $n$ 的多项式。选取首项系数 $2^n$ 的规范化，就得到下面的 $H_n$。
\ans{x=a_0\sum_{k\ge0}\left[\prod_{j=0}^{k-1}\dfrac{4j-\lambda}{(2j+2)(2j+1)}\right]t^{2k}
+a_1\sum_{k\ge0}\left[\prod_{j=0}^{k-1}\dfrac{4j+2-\lambda}{(2j+3)(2j+2)}\right]t^{2k+1},\\
\lambda=2n:\quad H_n(t)=n!\sum_{j=0}^{\lfloor n/2\rfloor}\dfrac{(-1)^j(2t)^{n-2j}}{j!(n-2j)!}.}
$H_n$ 的首项系数为 $2^n\ne0$；代入递推即可核验各系数满足原方程。$n=0$ 时为非零常数。
\problem{3.5.4}{87--88}{设复向量 $z=x+\ii y$ 为 $z'=A(t)z+f_R(t)+\ii f_I(t)$ 的复值解，其中系数为实。证明实部和虚部分别满足相应实方程。}
\begin{proof}
由 $A$ 为实矩阵，
\[
 x'+\ii y'=Ax+f_R+\ii(Ay+f_I).
\]
复数相等当且仅当实部、虚部分别相等，故 $x'=Ax+f_R$，$y'=Ay+f_I$。此结论依赖 $A$ 为实矩阵；若 $A$ 本身有虚部，则实虚部会耦合，不能直接套用。
逐个分量看，第 $j$ 个方程为 $x_j'+\ii y_j'=\sum_k a_{jk}x_k+(f_R)_j+\ii[\sum_k a_{jk}y_k+(f_I)_j]$。每个括号都是实数，因此实部和虚部分别给出两个实向量方程，不需要把实部与虚部的初值额外混合。
\end{proof}
